\documentclass[../../syllabus_1.tex]{subfiles}
\begin{document}

\chapter{Angles}
\section{Définition}
Un angle est une région du plan délimitée par deux demi-droites issues d'un point.

Un angle est saillant lorsqu'il peut etre contenu dans un angle plat.
\begin{center}
    \begin{tikzpicture}[baseline=(current bounding box.north)]
        \tkzDefPoint(0,0){A}
        \tkzDefPoint(30:2){B}
        \tkzDefPoint(2,0){C}
        \tkzFillAngle[blue!30](C,A,B)
        \tkzDrawLines[add=0 and 0.2](A,B A,C)
        \tkzDrawPoints(A,B,C)
    \end{tikzpicture}
    \hfil
    \begin{tikzpicture}[baseline=(current bounding box.north)]
        \tkzDefPoint(0,0){A}
        \tkzDefPoint(30:2){B}
        \tkzDefPoint(2,0){C}
        \tkzFillAngle[blue!30](B,A,C)
        \tkzDrawLines[add=0 and 0.2](A,B A,C)
        \tkzDrawPoints(C,B,A)
    \end{tikzpicture}
\end{center}

Le point dont sont issues les deux demi-droites est appelé sommet de l'angle et es demi-droites sont appelées cotés de l'angle.
\begin{center}
    \begin{tikzpicture}[baseline=(current bounding box.north)]
        \tkzDefPoint(0,0){A}
        \tkzDefPoint(30:2){B}
        \tkzDefPoint(2,0){C}
        \tkzDrawLines[add=0 and 0.2](A,B A,C)
        \tkzDrawPoints(A,B,C)
        \node (sommet) at (-1,-1){Sommet};
        \draw[->] (A) to[bend right] (sommet.north);
        \node (cotes) at (15:5){Cotés de l'angle};
        \draw [->] ($(A)!.5!(B)$) to[bend left] (cotes.north);
        \draw [->] ($(A)!.5!(C)$) to[bend right] (cotes.south);
    \end{tikzpicture}
\end{center}
\end{document}